% ============================================================================
% Fichier  : annexes/Ann_A_glossaire.tex
% Version  : 1.0 -- novembre 2026
% ============================================================================

\chapter{Glossaire}
\label{ann:glossaire}

Ce glossaire reprend, sous une forme abrégée, les définitions données dans
les chapitres. Les définitions tirées de l'article~5 de la loi n°~2024/017
sont paraphrasées~; le texte exact figure dans le chapitre indiqué.

\begin{description}
    \item[Anonymisation] Modification ou suppression d'informations pour
          rendre les données non attribuables à des personnes~; elle doit
          être prouvée (Ch.~\ref{chap:donnees}).
    \item[Autorisation préalable] Acte de l'Autorité auquel la loi soumet
          tout traitement (art.~19) et tout transfert (art.~32)
          (Ch.~\ref{chap:loi}).
    \item[Autorité de protection] Organisme public indépendant chargé de
          veiller à l'application de la loi, créé et organisé par décret
          (art.~53) (Ch.~\ref{chap:loi}).
    \item[Clauses contractuelles types] Clauses élaborées et publiées par
          l'Autorité pour encadrer les transferts (Ch.~\ref{chap:transferts}).
    \item[Consentement] Manifestation de volonté expresse, non équivoque, non
          contrainte, libre et éclairée~; préalable, spécifique et univoque
          selon l'article~9 (Ch.~\ref{chap:principes}).
    \item[Constatation] Fait vérifiable, rattaché à une pièce
          (Ch.~\ref{chap:prologue}).
    \item[Critère de substitution] Critère déclaré dans le plan d'audit
          quand le texte applicable n'est pas publié (Ch.~\ref{chap:droits}).
    \item[Décision automatisée] Décision fondée exclusivement sur un
          traitement automatisé (art.~44) (Ch.~\ref{chap:profilage}).
    \item[Découverte des données] Recherche, dans les supports réels, des
          données personnelles que le registre ignore
          (Ch.~\ref{chap:cartographie}).
    \item[Délégué à la protection des données] Fonction de bonne pratique,
          non prévue par la loi, chargée du registre, du conseil et des
          demandes (Ch.~\ref{chap:prologue}).
    \item[Destinataire] Personne habilitée à recevoir une communication des
          données (Ch.~\ref{chap:donnees}).
    \item[Donnée à caractère personnel] Information permettant d'identifier
          une personne directement ou indirectement (Ch.~\ref{chap:donnees}).
    \item[Donnée sensible] Donnée relative notamment aux opinions, aux
          transactions bancaires, à l'origine ethnique, linguistique ou
          régionale, à la vie sexuelle, à la génétique, à la biométrie, à la
          santé, aux poursuites et sanctions (Ch.~\ref{chap:donnees}).
    \item[Entropie ($H$)] Dispersion des cotes d'un ensemble d'exigences, en
          bits (Ch.~\ref{chap:audit}).
    \item[Évaluation d'impact] Procédure qui analyse la probabilité et la
          gravité des risques pour les droits et libertés (art.~33)
          (Ch.~\ref{chap:risque}).
    \item[Fichier] Ensemble structuré de données personnelles accessibles
          selon des critères (Ch.~\ref{chap:donnees}).
    \item[Finalité] But d'un traitement~; elle le définit et le limite
          (art.~11) (Ch.~\ref{chap:principes}).
    \item[Interconnexion] Mise en relation de données traitées pour une
          finalité avec d'autres données (art.~36)
          (Ch.~\ref{chap:cartographie}).
    \item[\NEK{} de critère] Incertitude qui tient à une règle non fixée
          (texte non publié, lecture non tranchée) (Ch.~\ref{chap:audit}).
    \item[\NEK{} de preuve] Incertitude qui tient à une pièce manquante
          (Ch.~\ref{chap:audit}).
    \item[\OUI{}, \NON{}, \NEK{}] Cotes d'une exigence~: conformité établie,
          écart établi, connaissance insuffisante (Ch.~\ref{chap:prologue}).
    \item[Papier de travail] Document qui trace un constat~: exigence,
          critère, test, preuve, cote, auteur (Ch.~\ref{chap:audit}).
    \item[Personne concernée] Personne dont les données sont traitées
          (Ch.~\ref{chap:donnees}).
    \item[Portabilité] Droit de recevoir gratuitement ses données dans un
          format structuré et de les faire transmettre (art.~43)
          (Ch.~\ref{chap:droits}).
    \item[Profilage] Traitement automatisé qui évalue des aspects personnels
          (santé, préférences, localisation, situation économique)
          (Ch.~\ref{chap:profilage}).
    \item[Pseudonymisation] Remplacement des identifiants directs par des
          identifiants artificiels~; les données restent personnelles
          (Ch.~\ref{chap:donnees}).
    \item[Registre des activités de traitement] Registre exigé par
          l'article~29 (Ch.~\ref{chap:cartographie}).
    \item[Responsable du traitement] Personne qui détermine les finalités et
          les moyens d'un traitement (Ch.~\ref{chap:donnees}).
    \item[Responsables conjoints] Responsables qui déterminent ensemble
          finalités et moyens~; convention publique (art.~31)
          (Ch.~\ref{chap:soustraitants}).
    \item[Sous-traitant] Personne qui traite des données pour le compte et
          sur instruction du responsable (Ch.~\ref{chap:donnees}).
    \item[Taux de conformité ($\mu$)] Part des exigences cotées \OUI{}
          (Ch.~\ref{chap:audit}).
    \item[Taux d'incertitude ($U$)] Part des exigences cotées \NEK{}~;
          décomposable en $U_p$ et $U_c$ (Ch.~\ref{chap:recherche}).
    \item[Tension normative] Situation où l'application simultanée de deux
          normes demande une interprétation, entre la compatibilité et
          l'incompatibilité (Ch.~\ref{chap:recherche}).
    \item[Traitement] Toute opération appliquée à des données personnelles,
          de la collecte à la destruction (Ch.~\ref{chap:donnees}).
    \item[Transfert international] Acheminement de données au-delà des
          frontières (Ch.~\ref{chap:transferts}).
    \item[Violation de données] Atteinte à la sécurité entraînant
          destruction, perte, altération, divulgation ou accès non autorisé
          (Ch.~\ref{chap:violations}).
\end{description}
